\pagenumbering{arabic}
\chapter{INTRODUCTION}
\section{Background}
Regular exercise contributes to a healthy lifestyle, but beginners often struggle to
maintain the correct form and keep an accurate record of their progress. Repeated
movements performed with poor posture can make a workout less effective and may also
increase the risk of injury. This project responds to that practical problem by
developing a \textbf{smart gym assistant} that uses \textbf{motion sensors} and
\textbf{machine learning} to recognize exercises, count repetitions, and identify
posture errors during exercise.


\section{Problem Definition}
Many gym users perform exercises without a trainer continuously checking their form,
and repetitions are often recorded manually. As a result, posture mistakes may go
unnoticed and workout records may be incomplete or inaccurate. The project therefore
aims to provide an accessible \textbf{AI-based assistant} that can track exercises
and give posture-correction feedback using affordable \textbf{embedded hardware}.


\section{Objectives}
\begin{itemize}
  \item To develop a smart gym assistant that monitors and analyzes exercise movements in real time.
  \item To recognize supported exercises, count repetitions, identify incorrect posture, and provide feedback that helps users improve their form and reduce the risk of injury.
\end{itemize}

\section{Scope of Project}
\begin{itemize}
  \item The current implementation focuses on squats and bicep curls using two \textbf{MPU6050 motion sensors} and a \textbf{KNN model}.
  \item The system identifies the exercise, detects good or bad posture, counts repetitions, and provides feedback through an \textbf{OLED}, \textbf{buzzer}, and \textbf{web interface}.
  \item The same data and model pipeline can later be extended to exercises such as push-ups by collecting and labeling additional training data.
\end{itemize}
